
Large language models exhibit two reliability concerns: generating factual errors and succumbing to adversarial perturbations. This paper proposes and tests the hypothesis that calibration quality---the alignment between model confidence and accuracy---serves as a shared internal signal enabling both factuality error detection and adversarial robustness. We develop a methodology evaluating open-weight LLMs on TruthfulQA (factuality), TextFooler attacks (robustness), and Expected Calibration Error (calibration), then analyzing whether ECE mediates the correlation between factuality and robustness metrics. We present a validated evaluation pipeline and analysis framework. A proof-of-concept evaluation on three models (FLAN-T5 base, FLAN-T5 large, Phi-2) demonstrates pipeline functionality. However, Attack Success Rate (ASR) values were simulated rather than measured, and only 3 of 12 planned models were evaluated, rendering all correlation findings artifacts of the simulation. This paper documents our methodology as a foundation for future work connecting these three reliability dimensions.

